\section{USAMO 1986}
        \begin{enumerate}
        	\item (\textit{USAMO 1986}) $(\text{a})$ Does there exist 14 consecutive positive integers each of which is divisible by one or more primes $p$ from the interval $2\le p \le 11$?


 $(\text{b})$ Does there exist 21 consecutive positive integers each of which is divisible by one or more primes $p$ from the interval $2\le p \le 13$?


 

	\item (\textit{USAMO 1986}) During a certain lecture, each of five mathematicians fell asleep exactly twice. For each pair of mathematicians, there was some moment when both were asleep simultaneously. Prove that, at some moment, three of them were sleeping simultaneously.


 

	\item (\textit{USAMO 1986}) What is the smallest integer $n$, greater than one, for which the root-mean-square of the first $n$ positive integers is an integer?


 $\mathbf{Note.}$ The root-mean-square of $n$ numbers $a_1, a_2, \cdots, a_n$ is defined to be

 \[\left[\frac{a_1^2 + a_2^2 + \cdots + a_n^2}n\right]^{1/2}\]


 

	\item (\textit{USAMO 1986}) Two distinct circles $K_1$ and $K_2$ are drawn in the plane. They intersect at points $A$ and $B$, where $AB$ is the diameter of $K_1$. A point $P$ on $K_2$ and inside $K_1$ is also given.


 Using only a "T-square" (i.e. an instrument which can produce a straight line joining two points and the perpendicular to a line through a point on or off the line), find a construction for two points $C$ and $D$ on $K_1$ such that $CD$ is perpendicular to $AB$ and $\angle CPD$ is a right angle.


 

	\item (\textit{USAMO 1986}) By a partition $\pi$ of an integer $n\ge 1$, we mean here a representation of $n$ as a sum of one or more positive integers where the summands must be put in nondecreasing order. (E.g., if $n=4$, then the partitions $\pi$ are $1+1+1+1$, $1+1+2$, $1+3, 2+2$, and $4$).


 For any partition $\pi$, define $A(\pi)$ to be the number of $1$'s which appear in $\pi$, and define $B(\pi)$ to be the number of distinct integers which appear in $\pi$. (E.g., if $n=13$ and $\pi$ is the partition $1+1+2+2+2+5$, then $A(\pi)=2$ and $B(\pi) = 3$).


 Prove that, for any fixed $n$, the sum of $A(\pi)$ over all partitions of $\pi$ of $n$ is equal to the sum of $B(\pi)$ over all partitions of $\pi$ of $n$.


 

\end{enumerate}